%%%%%%%%%%%%%%%%%%%%%%%%%%%%%%%%%%%%%%%%%%%%%%%%%%%%%%%%%%%%%%%%%%%%%%%%%%%%%%%
% CHAPTER 9: CONTEXT AS AN ENDOGENOUS VARIABLE
%%%%%%%%%%%%%%%%%%%%%%%%%%%%%%%%%%%%%%%%%%%%%%%%%%%%%%%%%%%%%%%%%%%%%%%%%%%%%%%
% Version: 2.0 (January 2026)
%
% Purpose: Introduces the 8 Ψ-dimensions of context and establishes context
%          as endogenous---shaped by the behavior it influences. Answers the
%          fourth CORE question: "When and how does context modulate behavior?"
%
% Primary Appendix: V - CORE-WHEN (The 8Ψ-Dimensions Framework)
% Secondary Appendices: AH (Temporal), AI (Spatial), AAA (WHO), AG (Complexity)
%
% Prerequisites: Ch. 5 (Complementarity), Ch. 8 (Mathematical Foundations)
% Leads to: Ch. 10 (Welfare/FEPSDE), Ch. 11 (Awareness)
%
% Chapter Type: A (CORE Chapter - answers WHEN question)
%%%%%%%%%%%%%%%%%%%%%%%%%%%%%%%%%%%%%%%%%%%%%%%%%%%%%%%%%%%%%%%%%%%%%%%%%%%%%%%

\chapter{Context as an Endogenous Variable}
\label{ch:context}

%%%%%%%%%%%%%%%%%%%%%%%%%%%%%%%%%%%%%%%%%%%%%%%%%%%%%%%%%%%%%%%%%%%%%%%%%%%%%%%
% CHAPTER SCOPE BOX
%%%%%%%%%%%%%%%%%%%%%%%%%%%%%%%%%%%%%%%%%%%%%%%%%%%%%%%%%%%%%%%%%%%%%%%%%%%%%%%
\begin{tcolorbox}[
    colback=purple!5!white,
    colframe=purple!75!black,
    title={\textbf{Chapter Scope: Ziel / In-Scope / Out-of-Scope}},
    fonttitle=\bfseries
]

\textbf{Ziel (Objective):}
Establish the formal foundation for the psi framework within the 10C CORE framework.

\vspace{0.3cm}
\textbf{In-Scope:}
\begin{itemize}[nosep]
    \item Core definitions and axioms for the psi framework
    \item Integration with other 10C dimensions
    \item Formal mathematical foundations
    \item Key theorems and their implications
\end{itemize}

\vspace{0.3cm}
\textbf{Out-of-Scope (delegiert an Appendices):}
\begin{itemize}[nosep]
    \item Detailed proofs and derivations $\rightarrow$ CORE Appendix
    \item Empirical calibration $\rightarrow$ Appendix BBB
    \item Domain-specific applications $\rightarrow$ Application chapters
\end{itemize}

\end{tcolorbox}



%%%%%%%%%%%%%%%%%%%%%%%%%%%%%%%%%%%%%%%%%%%%%%%%%%%%%%%%%%%%%%%%%%%%%%%%%%%%%%%
% QUICK REFERENCE BOX
%%%%%%%%%%%%%%%%%%%%%%%%%%%%%%%%%%%%%%%%%%%%%%%%%%%%%%%%%%%%%%%%%%%%%%%%%%%%%%%
\begin{tcolorbox}[
    colback=gray!5!white,
    colframe=gray!75!black,
    title={\textbf{Begriffe in diesem Kapitel}},
    fonttitle=\bfseries
]
\small
\begin{tabular}{@{}ll@{}}
\textbf{Key Concept} & \textbf{Definition} \\
\midrule
The Psi framework & See Section \ref{sec:ch9-intro} \\
\end{tabular}
\end{tcolorbox}



%%%%%%%%%%%%%%%%%%%%%%%%%%%%%%%%%%%%%%%%%%%%%%%%%%%%%%%%%%%%%%%%%%%%%%%%%%%%%%%
% INTUITION BOX
%%%%%%%%%%%%%%%%%%%%%%%%%%%%%%%%%%%%%%%%%%%%%%%%%%%%%%%%%%%%%%%%%%%%%%%%%%%%%%%
\begin{tcolorbox}[
    colback=blue!5!white,
    colframe=blue!75!black,
    title={\textbf{Intuition: A Motivating Example}},
    fonttitle=\bfseries
]
Consider \textbf{Anna}, a professional facing a decision that involves the psi framework.

She initially evaluates the choice using standard criteria, but something feels incomplete.
\textbf{Thomas}, her colleague, suggests considering additional dimensions beyond the obvious.

This chapter explores what Anna discovers when she applies the EBF framework---how
the psi framework interact with broader welfare considerations.

\medskip
\textit{By the end of this chapter, you will understand why Anna's initial analysis
was incomplete and how the EBF approach provides a more comprehensive view.}
\end{tcolorbox}



%==============================================================================
% DIE 9 FRAGEN NAVIGATION BOX
%==============================================================================
\BCMQuestionNav{9}

%==============================================================================
% 10C CORE CONNECTION BOX
%==============================================================================
\begin{tcolorbox}[
    colback=green!5!white,
    colframe=green!75!black,
    title={\textbf{10C CORE Connection: WHEN (Question 4)}},
    fonttitle=\bfseries
]
\textbf{The 4th CORE Question:} When and how does context modulate behavior?

\medskip
\begin{tabular}{ll}
\textbf{CORE Appendix:} & V (CORE-WHEN) \\
\textbf{Output:} & Context vector $\boldsymbol{\Psi} \in [0,1]^8$ \\
\textbf{Key Insight:} & Context is not background---it is the universal modulator \\
\textbf{Novel Contribution:} & Endogenous context: $\Psi_{t+1} = f(\Psi_t, a_t)$
\end{tabular}

\medskip
\textbf{The 8 Ψ-Dimensions:}
\begin{center}
\small
$\Psi_I$ (Institutional) $\cdot$ $\Psi_S$ (Social) $\cdot$ $\Psi_C$ (Cognitive) $\cdot$ $\Psi_K$ (Cultural)\\
$\Psi_E$ (Economic) $\cdot$ $\Psi_T$ (Temporal) $\cdot$ $\Psi_M$ (Material) $\cdot$ $\Psi_F$ (Physical)
\end{center}

\textit{Every parameter in EBF is context-dependent: $\gamma(\Psi)$, $\lambda(\Psi)$, $\theta(\Psi)$. This explains why ``one size fits all'' policies fail and treatment effects vary systematically.}
\end{tcolorbox}

%==============================================================================
% APPENDIX REFERENCES BOX
%==============================================================================
\begin{tcolorbox}[
    colback=yellow!10!white,
    colframe=orange!75!black,
    title={\textbf{Appendix References for This Chapter}},
    fonttitle=\bfseries
]
\begin{tabular}{lll}
\textbf{Appendix} & \textbf{Content} & \textbf{Section} \\
\midrule
V (CORE-WHEN) & 8 Ψ-dimensions, $\delta$-coefficients & All sections \\
AH (CONTEXT-TEMPORAL) & Time dimension details & \ref{sec:context-types} \\
AI (CONTEXT-SPATIAL) & Space dimension details & \ref{sec:context-spatial} \\
AAA (CORE-WHO) & Level salience $\alpha^L(\Psi)$ & \ref{sec:context-levels} \\
AG (DOMAIN-COMPLEXITY) & Emergent properties & \ref{sec:context-complexity} \\
G (REF-GLOSSARY) & Symbol definitions & All sections
\end{tabular}

\medskip
\textbf{Reading Guide:} This chapter develops the endogenous context concept;
Appendix WEN provides measurement protocols and calibration values.
\end{tcolorbox}

%==============================================================================
% CROSS-REFERENCE: EIT-EMP / EMERGE ALGORITHM
%==============================================================================
\begin{tcolorbox}[
    colback=blue!5!white,
    colframe=blue!75!black,
    title={\textbf{Cross-Reference: EMERGE Algorithm (Chapter 11, EIT-EMP)}},
    fonttitle=\bfseries
]
\textbf{WHEN $\rightarrow$ Universal Modulator:} Der Kontext-Vektor $\boldsymbol{\Psi} \in [0,1]^8$ moduliert \emph{alle} Parameter im EMERGE-Algorithmus (Chapter 11). Keine Intervention ist kontextunabhängig.

\medskip
\textbf{Ψ-Modulation der EMERGE-Faktoren:}
\begin{center}
\begin{tabular}{lll}
\textbf{EMERGE-Faktor} & \textbf{Ψ-Modulation} & \textbf{Beispiel} \\
\midrule
$f_\gamma$ (Crowding) & $\gamma(\Psi) \neq \gamma_0$ & Hohe $\Psi_S$ (sozial) $\Rightarrow$ stärkeres Crowding-Out \\
$f_S$ (Stage) & $\theta(\Psi) \neq \theta_0$ & Hohe $\Psi_E$ (ökonomisch) $\Rightarrow$ niedrigere Schwellen \\
$f_C$ (Complexity) & Feasibility $\sim \Psi$ & $\Psi_I$ (institutionell) bestimmt was möglich ist \\
$f_L$ (Scope) & $\alpha^L(\Psi)$ & Kontext bestimmt relevante Scope-Levels \\
\end{tabular}
\end{center}

\medskip
\textbf{Endogene Kontextdynamik:}
\[
\Psi_{t+1} = f(\Psi_t, a_t) \quad \Rightarrow \quad K^*_{t+1} = K^*(\Psi_{t+1}) \neq K^*_t
\]
Interventionen verändern den Kontext, was die optimale Portfoliogröße $K^*$ im Zeitverlauf verschiebt. Das EMERGE-Algorithmus muss daher dynamisch re-evaluiert werden (siehe Chapter 21: Dynamic Portfolio Evolution).

\medskip
\textbf{Die 8 Ψ-Dimensionen im EMERGE-Kontext:}
\begin{center}
\small
\begin{tabular}{llll}
$\Psi_I$ (Institutional) & $\Psi_S$ (Social) & $\Psi_C$ (Cognitive) & $\Psi_K$ (Cultural) \\
$\Psi_E$ (Economic) & $\Psi_T$ (Temporal) & $\Psi_M$ (Material) & $\Psi_F$ (Physical)
\end{tabular}
\end{center}

\medskip
\textbf{Workflow:} Chapter 9 (WHEN: $\Psi$) $\rightarrow$ moduliert alle 10C $\rightarrow$ Chapter 11 (EMERGE: $K^*(\Psi)$)

\medskip
\textit{Siehe Chapter 11, Section 11.5 für die vollständige EMERGE-Spezifikation.}
\end{tcolorbox}


%%%%%%%%%%%%%%%%%%%%%%%%%%%%%%%%%%%%%%%%%%%%%%%%%%%%%%%%%%%%%%%%%%%%%%%%%%%%%%%
%                          PART I: INTRODUCTION
%%%%%%%%%%%%%%%%%%%%%%%%%%%%%%%%%%%%%%%%%%%%%%%%%%%%%%%%%%%%%%%%%%%%%%%%%%%%%%%

\section{Introduction: Why Context Is Not Background}
\label{sec:context-intro}

%------------------------------------------------------------------------------

The preceding chapters treated context $\Psi$ as a parameter that shapes
the reference structure $C^*(\Psi)$ and the payoff landscape $\Pi(s, \sigma; \Psi)$.
But where does context come from? How does it change?

\begin{intuition}[The Feedback Loop]
Consider trust in a society:
\begin{itemize}[nosep]
  \item High trust ($\Psi_t$) $\Rightarrow$ Cooperation pays ($C^*$)
  \item Cooperation pays $\Rightarrow$ People cooperate ($a_t$)
  \item People cooperate $\Rightarrow$ Trust reinforced ($\Psi_{t+1}$)
\end{itemize}
But also:
\begin{itemize}[nosep]
  \item Low trust ($\Psi_t$) $\Rightarrow$ Defection pays ($C^*$)
  \item Defection pays $\Rightarrow$ People defect ($a_t$)
  \item People defect $\Rightarrow$ Trust eroded ($\Psi_{t+1}$)
\end{itemize}
Both are stable equilibria. Which one emerges depends on history.
Context is not exogenous---it is shaped by the behavior it influences.
\end{intuition}

%------------------------------------------------------------------------------
% CENTRAL QUESTION BOX
%------------------------------------------------------------------------------
\paragraph{The Central Question.}
This chapter addresses the 4th of the 9 BCM Questions:

\begin{center}
\fcolorbox{green!75!black}{green!5!white}{\parbox{0.85\textwidth}{
\textbf{Question 4 (WHEN):} When and how does context modulate behavior?
What are the systematic dimensions of context?
}}
\end{center}

\begin{definition}[Endogenous Context]
Context evolves through aggregate behavior:
\begin{equation}
\Psi_{t+1} = f(\Psi_t, a_t)
\label{eq:endogenous-context}
\end{equation}
Actions today shape the context tomorrow. This closes the feedback loop
that makes economic systems truly dynamic.
\end{definition}

%------------------------------------------------------------------------------
% CHAPTER OVERVIEW
%------------------------------------------------------------------------------
\paragraph{Chapter Overview.}
This chapter proceeds as follows:
\begin{itemize}
    \item Section \ref{sec:context-components}: The components of context
    \item Section \ref{sec:context-feedback}: The feedback loop and dynamics
    \item Section \ref{sec:context-types}: Different context types and speeds
    \item Section \ref{sec:context-tipping}: Tipping points and regime change
    \item Section \ref{sec:context-spatial}: Spatial context and agglomeration
    \item Section \ref{sec:context-complexity}: Complexity economics integration
    \item Section \ref{sec:context-levels}: Context at each analytical level
    \item Section \ref{sec:context-policy}: Policy implications
    \item Section \ref{sec:context-measurement}: Measuring context with NLP/LLMs
    \item Section \ref{sec:context-integration}: Integration with 10C framework
    \item Section \ref{sec:context-summary}: Summary and transition
\end{itemize}


%%%%%%%%%%%%%%%%%%%%%%%%%%%%%%%%%%%%%%%%%%%%%%%%%%%%%%%%%%%%%%%%%%%%%%%%%%%%%%%
%                     PART II: WHAT IS CONTEXT?
%%%%%%%%%%%%%%%%%%%%%%%%%%%%%%%%%%%%%%%%%%%%%%%%%%%%%%%%%%%%%%%%%%%%%%%%%%%%%%%

\section{The Components of Context}
\label{sec:context-components}

%------------------------------------------------------------------------------

\paragraph{The Components of $\Psi$.}
Context encompasses everything that shapes the payoff structure
but is not under any single agent's immediate control:

\begin{table}[h]
\centering
\begin{tabular}{lll}
\toprule
\textbf{Component} & \textbf{Examples} & \textbf{Time Scale} \\
\midrule
Technology & Production methods, digital infrastructure & Years to decades \\
Institutions & Laws, property rights, constitutions & Decades to centuries \\
Norms & Social expectations, fairness standards & Years to generations \\
Culture & Values, worldviews, identity categories & Generations \\
Physical environment & Climate, geography, infrastructure & Decades to centuries \\
\bottomrule
\end{tabular}
\caption{Context components and their characteristic time scales}
\label{tab:context-components}
\end{table}

\paragraph{Why Classical Economics Treats $\Psi$ as Exogenous.}
Arrow-Debreu and most microeconomic models assume:
\begin{equation}
\Psi_{t+1} = \Psi_t = \bar{\Psi}
\label{eq:exogenous-assumption}
\end{equation}

This simplification enables equilibrium analysis but rules out:
\begin{itemize}[nosep]
  \item Institutional change
  \item Technological progress
  \item Cultural evolution
  \item Path dependence
\end{itemize}

\begin{cross-reference}
See Appendix WEN (CORE-WHEN) for the complete specification of all 8 Ψ-dimensions
with measurement protocols and $\delta$-modulation coefficients.
\end{cross-reference}


%%%%%%%%%%%%%%%%%%%%%%%%%%%%%%%%%%%%%%%%%%%%%%%%%%%%%%%%%%%%%%%%%%%%%%%%%%%%%%%
%                     PART III: THE FEEDBACK LOOP
%%%%%%%%%%%%%%%%%%%%%%%%%%%%%%%%%%%%%%%%%%%%%%%%%%%%%%%%%%%%%%%%%%%%%%%%%%%%%%%

\section{The Feedback Loop}
\label{sec:context-feedback}

%------------------------------------------------------------------------------

\paragraph{The Circular Causation.}

\begin{definition}[Context Feedback Loop]
\begin{center}
\begin{tabular}{c}
$\Psi_t$ (Context) \\
$\downarrow$ \\
$C^*(\Psi_t)$ (Reference Structure) \\
$\downarrow$ \\
$a_t^*(C^*)$ (Optimal Actions) \\
$\downarrow$ \\
$\Psi_{t+1} = f(\Psi_t, a_t)$ (New Context) \\
$\downarrow$ \\
[cycle repeats]
\end{tabular}
\end{center}
\end{definition}

\subsection{A Simple Model of Context Dynamics}

\begin{definition}[Linear Context Update]
Near equilibrium, context dynamics can be approximated as:
\begin{equation}
\Psi_{t+1} = \Psi_t + \eta \cdot (a_t - a^*(\Psi_t))
\label{eq:context-update}
\end{equation}
where:
\begin{itemize}[nosep]
  \item $\eta > 0$: Adjustment speed (how fast context responds to behavior)
  \item $a_t$: Actual aggregate behavior at time $t$
  \item $a^*(\Psi_t)$: Equilibrium behavior given context $\Psi_t$
\end{itemize}
\end{definition}

\begin{intuition}[Context Dynamics]
If people behave \emph{as expected} ($a_t = a^*$), context doesn't change.
If people behave \emph{differently} ($a_t \neq a^*$), context shifts.

\textbf{Example: Norm Erosion.}
\begin{itemize}[nosep]
  \item If everyone follows the norm ($a_t = a^* = 0.05$ littering rate): Norm persists
  \item If deviation occurs ($a_t = 0.15 > a^*$): Norm weakens by $\Delta\Psi = -0.08$
  \item Weaker norm $\Rightarrow$ acceptable littering rises to 0.20 $\Rightarrow$ more littering
  \item After 5 periods: norm can collapse to $a^* = 0.40$ (8× baseline)
\end{itemize}
\end{intuition}


%%%%%%%%%%%%%%%%%%%%%%%%%%%%%%%%%%%%%%%%%%%%%%%%%%%%%%%%%%%%%%%%%%%%%%%%%%%%%%%
%                     PART IV: CONTEXT TYPES AND SPEEDS
%%%%%%%%%%%%%%%%%%%%%%%%%%%%%%%%%%%%%%%%%%%%%%%%%%%%%%%%%%%%%%%%%%%%%%%%%%%%%%%

\section{Different Types of Context: Speed of Change}
\label{sec:context-types}

%------------------------------------------------------------------------------

Not all context changes at the same speed.

\subsection{Fast-Moving Context: Technology ($\eta_{tech}$ large)}

\begin{itemize}[nosep]
  \item Can change in years (smartphones: 2007-2015)
  \item Driven by innovation, investment, adoption
  \item Individuals can influence through entrepreneurship
\end{itemize}

\subsection{Medium-Moving Context: Norms ($\eta_{norm}$ medium)}

\begin{itemize}[nosep]
  \item Changes over years to decades (smoking norms: 1970-2000)
  \item Driven by social learning, signaling, coordination
  \item Requires critical mass to shift
\end{itemize}

\subsection{Slow-Moving Context: Institutions ($\eta_{inst}$ small)}

\begin{itemize}[nosep]
  \item Changes over decades to centuries (property rights, democracy)
  \item Driven by political conflict, crises, external pressure
  \item Highly persistent---``sticky''
\end{itemize}

\subsection{Very Slow Context: Culture ($\eta_{culture}$ very small)}

\begin{itemize}[nosep]
  \item Changes over generations (religious values, family structures)
  \item Driven by socialization, education, demographic change
  \item Extremely persistent
\end{itemize}

\paragraph{Implication: Layered Dynamics.}
\begin{definition}[Multi-Speed Context]
\begin{equation}
\Psi = (\Psi_{tech}, \Psi_{norm}, \Psi_{inst}, \Psi_{culture})
\label{eq:layered-context}
\end{equation}
Each component has its own dynamics:
\begin{align}
\Psi_{tech, t+1} &= \Psi_{tech, t} + \eta_{tech} \cdot \Delta a_t \\
\Psi_{norm, t+1} &= \Psi_{norm, t} + \eta_{norm} \cdot \Delta a_t \\
\Psi_{inst, t+1} &= \Psi_{inst, t} + \eta_{inst} \cdot \Delta a_t \\
\Psi_{culture, t+1} &= \Psi_{culture, t} + \eta_{culture} \cdot \Delta a_t
\end{align}
Technology moves fast; culture moves slow. This creates ``cultural lag'' \citep{ogburn1922}.
\end{definition}

\begin{cross-reference}
See Appendix CTM (CONTEXT-TEMPORAL) for detailed temporal dynamics
and Appendix WEN (CORE-WHEN) for the complete speed hierarchy.
\end{cross-reference}


%%%%%%%%%%%%%%%%%%%%%%%%%%%%%%%%%%%%%%%%%%%%%%%%%%%%%%%%%%%%%%%%%%%%%%%%%%%%%%%
%                     PART V: TIPPING POINTS
%%%%%%%%%%%%%%%%%%%%%%%%%%%%%%%%%%%%%%%%%%%%%%%%%%%%%%%%%%%%%%%%%%%%%%%%%%%%%%%

\section{Tipping Points and Regime Change}
\label{sec:context-tipping}

%------------------------------------------------------------------------------

\subsection{Self-Reinforcing vs. Self-Correcting Dynamics}

\begin{definition}[Feedback Types]
\textbf{Self-Reinforcing (Positive Feedback):}
\begin{equation}
\frac{\partial^2 \Psi_{t+1}}{\partial \Psi_t \partial a_t} > 0
\end{equation}
Example: Network effects---more users make platform more valuable.

\textbf{Self-Correcting (Negative Feedback):}
\begin{equation}
\frac{\partial^2 \Psi_{t+1}}{\partial \Psi_t \partial a_t} < 0
\end{equation}
Example: Price mechanisms---high demand raises prices, reducing demand.
\end{definition}

\paragraph{Which Dominates?}
The stability of context depends on whether self-reinforcing or
self-correcting dynamics dominate:
\begin{itemize}[nosep]
  \item Self-correcting dominant: Single stable equilibrium
  \item Self-reinforcing dominant: Multiple equilibria, path dependence
  \item Balance: Tipping points, regime shifts
\end{itemize}

\subsection{Nonlinear Dynamics and Thresholds}

\begin{definition}[Nonlinear Context Update]
\begin{equation}
\Psi_{t+1} = \Psi_t + \eta \cdot g(\Psi_t) \cdot (a_t - a^*(\Psi_t))
\label{eq:nonlinear-context}
\end{equation}
where $g(\Psi)$ varies with context:
\begin{itemize}[nosep]
  \item $g(\Psi) \approx 0$ when context is ``locked in''
  \item $g(\Psi)$ large when context is ``fluid''
\end{itemize}
\end{definition}

\begin{intuition}[Revolution as Tipping Point]
\begin{itemize}[nosep]
  \item Normal times: $g(\Psi) \approx 0$, institutions resist change
  \item Crisis: $g(\Psi)$ suddenly large, small actions have big effects
  \item New equilibrium: $g(\Psi) \approx 0$ again, new institutions lock in
\end{itemize}
\end{intuition}

\subsection{Historical Examples}

\paragraph{The Fall of Communism (1989).}
\begin{itemize}[nosep]
  \item For decades: $g(\Psi) \approx 0$, system appeared stable
  \item Gorbachev reforms: $g(\Psi)$ increased, system became fluid
  \item Small protests in Leipzig $\Rightarrow$ massive cascade
  \item Within months: Complete regime change across Eastern Europe
\end{itemize}

\paragraph{The Smartphone Revolution (2007-2015).}
\begin{itemize}[nosep]
  \item 2007: iPhone launches---$\Psi_{tech}$ begins shifting
  \item 2010: Critical mass reached---app ecosystem explodes
  \item 2015: Smartphones ubiquitous---new equilibrium
\end{itemize}

\paragraph{The Decline of Smoking (1965-2020).}
Norm change followed S-curve over 55 years:
\begin{equation}
\Psi_{t+1} = \Psi_t + \eta \cdot \Psi_t (1 - \Psi_t) \cdot \Delta a_t
\label{eq:s-curve}
\end{equation}
Early stage slow (few adopters), middle stage fast (critical mass), late stage slow (saturation).


%%%%%%%%%%%%%%%%%%%%%%%%%%%%%%%%%%%%%%%%%%%%%%%%%%%%%%%%%%%%%%%%%%%%%%%%%%%%%%%
%                     PART VI: SPATIAL CONTEXT
%%%%%%%%%%%%%%%%%%%%%%%%%%%%%%%%%%%%%%%%%%%%%%%%%%%%%%%%%%%%%%%%%%%%%%%%%%%%%%%

\section{Spatial Context: Geography as Endogenous}
\label{sec:context-spatial}

%------------------------------------------------------------------------------

Physical location constitutes a fundamental dimension of context ($\Psi_P$).
The New Economic Geography \citep{krugman1991} showed how spatial equilibria
emerge from the interaction of increasing returns, transport costs, and factor mobility.

\subsection{Agglomeration as Endogenous Context}

\citet{marshall1890} identified three sources of agglomeration benefits:
\begin{enumerate}[nosep]
  \item \textbf{Labor pooling:} Thick markets reduce matching frictions
  \item \textbf{Input sharing:} Specialized suppliers viable at scale
  \item \textbf{Knowledge spillovers:} Ideas flow locally
\end{enumerate}

In our notation, these are spatial complementarities:
\begin{equation}
C_{ij}^{spatial} > 0 \quad \text{(co-location is complementary)}
\label{eq:spatial-complementarity}
\end{equation}

\subsection{Core-Periphery Dynamics}

\begin{itemize}[nosep]
  \item Multiple equilibria: Both dispersed and agglomerated configurations are stable
  \item Path dependence: Small initial advantages can lock in spatial structure
  \item Discontinuous transitions: Gradual parameter change can trigger sudden concentration
\end{itemize}

\begin{cross-reference}
See Appendix CTS (CONTEXT-SPATIAL) for detailed spatial dynamics
and agglomeration economics integration.
\end{cross-reference}


%%%%%%%%%%%%%%%%%%%%%%%%%%%%%%%%%%%%%%%%%%%%%%%%%%%%%%%%%%%%%%%%%%%%%%%%%%%%%%%
%                     PART VII: COMPLEXITY ECONOMICS
%%%%%%%%%%%%%%%%%%%%%%%%%%%%%%%%%%%%%%%%%%%%%%%%%%%%%%%%%%%%%%%%%%%%%%%%%%%%%%%

\section{Complexity Economics: Emergence and Non-Equilibrium}
\label{sec:context-complexity}

%------------------------------------------------------------------------------

Standard economics assumes equilibrium exists, is unique, and is reached quickly.
Complexity economics \citep{arthur2015} challenges all three assumptions.

\subsection{Economy as Complex Adaptive System}

The Santa Fe approach \citep{arthur1997} characterizes the economy as:
\begin{itemize}[nosep]
  \item \textbf{Heterogeneous agents:} Not identical, not representative
  \item \textbf{Local interaction:} Agents interact with neighbors, not all
  \item \textbf{Bounded rationality:} $\Psi_C < \infty$
  \item \textbf{Adaptation:} Agents learn, strategies evolve
  \item \textbf{Emergence:} Macro patterns not derivable from micro rules
\end{itemize}

\subsection{Increasing Returns and Lock-In}

\citet{arthur1989,arthur1994} showed that with increasing returns, small early
events can have large long-run consequences:
\begin{equation}
C^*_{selected} \in \{C^*_1, C^*_2, \ldots\} \quad \text{history selects which}
\label{eq:lock-in}
\end{equation}

\subsection{Emergence and Segregation}

\citet{schelling1971,schelling1978} demonstrated that macro patterns emerge
from micro behavior in unexpected ways. His segregation model shows that
mild individual preferences produce complete segregation:
\begin{equation}
C^*_{macro} \neq \sum_i C^*_i \quad \text{(emergent properties)}
\label{eq:emergence}
\end{equation}

\begin{cross-reference}
See Appendix SHA (DOMAIN-COMPLEXITY) for detailed complexity economics integration
and agent-based modeling approaches.
\end{cross-reference}


%%%%%%%%%%%%%%%%%%%%%%%%%%%%%%%%%%%%%%%%%%%%%%%%%%%%%%%%%%%%%%%%%%%%%%%%%%%%%%%
%                     PART VIII: CONTEXT AT EACH LEVEL
%%%%%%%%%%%%%%%%%%%%%%%%%%%%%%%%%%%%%%%%%%%%%%%%%%%%%%%%%%%%%%%%%%%%%%%%%%%%%%%

\section{Context Dynamics at Each Level}
\label{sec:context-levels}

%------------------------------------------------------------------------------

Context operates differently at each analytical level.

\subsection{Individual Level: Personal Context}

Components: Habits, routines, physical environment, social network, skills.
\begin{equation}
\Psi_{personal, t+1} = \Psi_{personal, t} + \eta_{habit} \cdot (a_t - a^*)
\end{equation}

\subsection{Household Level: Family Context}

Components: Role expectations, communication patterns, resource allocation rules.
\begin{equation}
\Psi_{family, t+1} = \Psi_{family, t} + \sum_i w_i \cdot (a_{i,t} - a_i^*)
\end{equation}

\subsection{Organizational Level: Corporate Context}

Components: Corporate culture, formal structure, informal power networks.
\begin{equation}
\Psi_{org, t+1} = \Psi_{org, t} + \eta_{org}(leadership_t) \cdot \Delta strategy_t
\end{equation}

\subsection{National Level: Institutional Context}

Components: Constitutional framework, legal system, economic institutions.
\begin{equation}
\Psi_{national, t+1} = \begin{cases}
\Psi_{national, t} & \text{normal times} \\
\Psi_{national, t} + \Delta_{crisis} & \text{critical juncture}
\end{cases}
\end{equation}

\subsection{Global Level: International Context}

Components: International organizations, trade regimes, global norms.
\begin{equation}
\Psi_{global, t+1} = \Psi_{global, t} + \eta_{hegemony} \cdot a_{hegemon} + \eta_{crisis} \cdot crisis_t
\end{equation}

\begin{cross-reference}
See Appendix WHO (CORE-WHO) for the complete level hierarchy
and aggregation weights $\alpha^L(\Psi)$.
\end{cross-reference}


%%%%%%%%%%%%%%%%%%%%%%%%%%%%%%%%%%%%%%%%%%%%%%%%%%%%%%%%%%%%%%%%%%%%%%%%%%%%%%%
%                     PART IX: POLICY IMPLICATIONS
%%%%%%%%%%%%%%%%%%%%%%%%%%%%%%%%%%%%%%%%%%%%%%%%%%%%%%%%%%%%%%%%%%%%%%%%%%%%%%%

\section{Policy Implications: Changing Context vs. Changing Behavior}
\label{sec:context-policy}

%------------------------------------------------------------------------------

\subsection{The Traditional Approach: Change Behavior}

Most policies target behavior directly:
\begin{itemize}[nosep]
  \item Taxes and subsidies (change incentives)
  \item Regulations (mandate or prohibit behavior)
  \item Information campaigns (change beliefs)
\end{itemize}

\textbf{Problem:} If context ($\Psi$) is unchanged, behavior reverts when policy is removed.

\subsection{The Context Approach: Change $\Psi$ Directly}

Alternative: target context itself:
\begin{itemize}[nosep]
  \item \textbf{Technology policy:} R\&D funding, infrastructure investment
  \item \textbf{Norm entrepreneurship:} Role models, media campaigns
  \item \textbf{Institutional reform:} Constitutional change, legal reform
  \item \textbf{Cultural policy:} Education, national narratives
\end{itemize}

\textbf{Advantage:} Changed context sustains changed behavior.

\subsection{When to Target Context vs. Behavior}

\begin{table}[h]
\centering
\begin{tabular}{lll}
\toprule
\textbf{Situation} & \textbf{Target} & \textbf{Example} \\
\midrule
Quick fix needed & Behavior & Tax on plastic bags \\
Sustainable change & Context & Norm against single-use plastic \\
Individual problem & Behavior & Fine for littering \\
Systemic problem & Context & Redesign waste system \\
\bottomrule
\end{tabular}
\caption{When to target context vs. behavior}
\label{tab:policy-targeting}
\end{table}

\begin{intuition}[The Sequencing Problem]
Optimal policy often sequences both:
\begin{enumerate}[nosep]
  \item Start behavior change (regulations, incentives)
  \item Behavior change shifts norms
  \item Norm change enables institutional reform
  \item New institutions sustain new behavior
  \item Remove original policy---new equilibrium is self-sustaining
\end{enumerate}
Example: Seat belt laws $\rightarrow$ wearing becomes normal $\rightarrow$ law rarely enforced.
\end{intuition}


%%%%%%%%%%%%%%%%%%%%%%%%%%%%%%%%%%%%%%%%%%%%%%%%%%%%%%%%%%%%%%%%%%%%%%%%%%%%%%%
%                     PART X: MEASURING CONTEXT
%%%%%%%%%%%%%%%%%%%%%%%%%%%%%%%%%%%%%%%%%%%%%%%%%%%%%%%%%%%%%%%%%%%%%%%%%%%%%%%

\section{Measuring Context: The NLP/LLM Breakthrough}
\label{sec:context-measurement}

%------------------------------------------------------------------------------

Context $\Psi$ includes norms, trust, culture---variables that traditionally
resisted quantification. Modern NLP has transformed this.

\subsection{The NLP Revolution (2020-2026)}

\textbf{1. Norm Extraction from Text.}
Large corpora contain implicit norms. Modern NLP can extract:
\begin{itemize}[nosep]
  \item What behaviors are praised or condemned
  \item How justifications shift over time
  \item When norm change reaches tipping points
\end{itemize}

\textbf{2. Trust from Communication Patterns.}
Trust manifests in language: cooperative vs. defensive framing,
future-oriented vs. present-oriented discourse.

\textbf{3. Institutional Quality from Legal Text.}
Laws and regulations are text. Modern NLP can compare legal frameworks,
track regulatory change, identify gaps between law and enforcement.

\subsection{LLM-Powered Context Analysis}

\begin{intuition}[From Qualitative to Quantitative $\Psi$]
Example prompt:
\begin{quote}
``Analyze these 1000 policy documents from Germany (2000-2025).
Extract: (1) implicit assumptions about labor market flexibility,
(2) attitudes toward risk-sharing between firms and workers.
Score each dimension 0-10 and track trends.''
\end{quote}
This produces quantified context measures in hours rather than years.
\end{intuition}

\subsection{The Calibration Gap}

\begin{definition}[Two Tasks]
\begin{center}
\begin{tabular}{lll}
\toprule
\textbf{Task} & \textbf{LLM Capability} & \textbf{Requirement} \\
\midrule
Measure $\Psi$ & Strong & Text analysis at scale \\
Predict $C^*(\Psi)$ & Weak & Calibration on behavioral data \\
\bottomrule
\end{tabular}
\end{center}
LLMs measure context but cannot predict behavioral responses without calibration.
\end{definition}

\begin{cross-reference}
See Appendix WEN (CORE-WHEN), Section 8 for complete measurement protocols
and LLM-MC estimation methodology.
\end{cross-reference}


%%%%%%%%%%%%%%%%%%%%%%%%%%%%%%%%%%%%%%%%%%%%%%%%%%%%%%%%%%%%%%%%%%%%%%%%%%%%%%%
%                     PART XI: 10C INTEGRATION
%%%%%%%%%%%%%%%%%%%%%%%%%%%%%%%%%%%%%%%%%%%%%%%%%%%%%%%%%%%%%%%%%%%%%%%%%%%%%%%

\section{Integration with the CORE Framework}
\label{sec:context-integration}

%------------------------------------------------------------------------------

\subsection{Connection to CORE-WHO (Levels)}

\begin{cross-reference}
Context modulates which level is salient:
\begin{itemize}[nosep]
  \item Level salience function $\alpha^L(\Psi)$
  \item Crisis context: Lower levels become salient (self-preservation)
  \item Collective context: Higher levels become salient (group identity)
\end{itemize}
See Appendix WHO (CORE-WHO) for the complete level specification.
\end{cross-reference}

\subsection{Connection to CORE-WHAT (FEPSDE)}

\begin{cross-reference}
Context shifts dimension weights:
\begin{itemize}[nosep]
  \item Poverty: $U^F$ dominates
  \item Security: $U^P$ dominates
  \item Prosperity: $U^E$, $U^S$ gain weight
\end{itemize}
See Chapter 10 and Appendix WAT (CORE-WHAT) for the utility specification.
\end{cross-reference}

\subsection{Connection to CORE-HOW (Complementarity)}

\begin{cross-reference}
Context modulates complementarity:
\begin{equation}
\gamma(\Psi) = \gamma_0 + \sum_k \delta_k \cdot \Psi_k
\end{equation}
See Chapter 5 and Appendix HOW (CORE-HOW) for the complementarity specification.
\end{cross-reference}

\subsection{Connection to CORE-AWARE and CORE-READY}

\begin{cross-reference}
Context affects both awareness and willingness:
\begin{itemize}[nosep]
  \item Salience $s^d(\Psi)$: Context affects what's noticed
  \item Threshold $\theta(\Psi)$: Context affects action barriers
\end{itemize}
See Chapters 11-12 and Appendices AU, AV.
\end{cross-reference}

\subsection{Connection to CORE-STAGE (BCJ)}

\begin{cross-reference}
Context affects stage transitions:
\begin{itemize}[nosep]
  \item $\Delta\Psi$ can trigger stage transitions
  \item Context stability enables habit formation ($S_6$)
  \item Context change can cause relapse
\end{itemize}
See Chapter 13 and Appendix STA (CORE-STAGE).
\end{cross-reference}

%------------------------------------------------------------------------------
\subsection{The 10C Integration Table}

\begin{table}[h]
\centering
\begin{tabular}{lll}
\toprule
\textbf{CORE} & \textbf{Question} & \textbf{Context Role} \\
\midrule
WHO (AAA) & Who has utility? & Level salience $\alpha^L(\Psi)$ \\
WHAT (C) & What is utility? & Dimension weights $\nu_d(\Psi)$ \\
HOW (B) & How interact? & $\gamma(\Psi)$ complementarity \\
\textbf{WHEN (V)} & \textbf{When context?} & \textbf{8 $\Psi$-dimensions} \\
WHERE (BBB) & Where parameters? & $\Theta(\Psi)$ context-dependent \\
AWARE (AU) & How aware? & Salience $s^d(\Psi)$ \\
READY (AV) & How willing? & Threshold $\theta(\Psi)$ \\
STAGE (AW) & Where in journey? & Stage transitions via $\Delta\Psi$ \\
\bottomrule
\end{tabular}
\caption{Integration of WHEN (Context) with the 10C framework}
\label{tab:10c-context-integration}
\end{table}


%%%%%%%%%%%%%%%%%%%%%%%%%%%%%%%%%%%%%%%%%%%%%%%%%%%%%%%%%%%%%%%%%%%%%%%%%%%%%%%
%                     PART XII: THE COMPLETE SYSTEM
%%%%%%%%%%%%%%%%%%%%%%%%%%%%%%%%%%%%%%%%%%%%%%%%%%%%%%%%%%%%%%%%%%%%%%%%%%%%%%%

\section{The Complete Dynamic System}
\label{sec:context-system}

%------------------------------------------------------------------------------

Combining context dynamics with coherence dynamics:

\begin{definition}[Complete Dynamic System]
\begin{align}
C_{t+1} &= \Phi(C_t, \Psi_t) && \text{(Behavior adapts to context)} \\
\Psi_{t+1} &= f(\Psi_t, a_t(C_t)) && \text{(Context adapts to behavior)} \\
K_t &= 1 - \|C_t - C^*(\Psi_t)\| / \|C^*(\Psi_t)\| && \text{(Coherence measure)} \\
Q_t &= \Pi(C_t; \Psi_t) && \text{(Quality measure)}
\end{align}
\end{definition}

\paragraph{Steady State.}
At steady state:
\begin{itemize}[nosep]
  \item $C = C^*(\Psi)$: Behavior matches beliefs
  \item $\Psi_{t+1} = \Psi_t$: Context is stable
  \item $K = 1$: Full coherence
  \item $Q = Q(C^*; \Psi)$: Quality determined by which equilibrium
\end{itemize}

\paragraph{Transitions.}
During transitions:
\begin{itemize}[nosep]
  \item $C \neq C^*$: Behavior and beliefs diverge
  \item $\Psi$ changing: Context in flux
  \item $K < 1$: Incoherence
  \item $Q$ may rise or fall depending on direction
\end{itemize}


%%%%%%%%%%%%%%%%%%%%%%%%%%%%%%%%%%%%%%%%%%%%%%%%%%%%%%%%%%%%%%%%%%%%%%%%%%%%%%%
%                          PART XIII: SUMMARY
%%%%%%%%%%%%%%%%%%%%%%%%%%%%%%%%%%%%%%%%%%%%%%%%%%%%%%%%%%%%%%%%%%%%%%%%%%%%%%%

\section{Summary}
\label{sec:context-summary}

%------------------------------------------------------------------------------

This chapter has established:

\begin{enumerate}
    \item \textbf{Endogenous Context:} $\Psi_{t+1} = f(\Psi_t, a_t)$---context is shaped
          by the behavior it influences, creating feedback loops.

    \item \textbf{Multi-Speed Dynamics:} Technology (fast), norms (medium), institutions (slow),
          culture (very slow)---each with characteristic $\eta$ values.

    \item \textbf{Tipping Points:} Nonlinear dynamics create thresholds where small
          changes trigger regime shifts.

    \item \textbf{Spatial Context:} Geography as endogenous through agglomeration
          and path dependence.

    \item \textbf{Complexity Integration:} Emergence, lock-in, and non-equilibrium
          dynamics from complexity economics.

    \item \textbf{Universal Modulator:} Every EBF parameter is context-dependent:
          $\gamma(\Psi)$, $\lambda(\Psi)$, $\theta(\Psi)$, $\alpha^L(\Psi)$.
\end{enumerate}

%------------------------------------------------------------------------------
\paragraph{Formal Details.}
The complete formal treatment appears in Appendix~V (CORE-WHEN), which contains:
\begin{itemize}[nosep]
    \item Complete 8 Ψ-dimension specification
    \item $\delta$-modulation coefficient matrices
    \item Context measurement protocols
    \item LLM-MC estimation methodology
    \item 80+ scientific references
\end{itemize}

%------------------------------------------------------------------------------
\paragraph{What Comes Next.}
Chapter 10 introduces \textbf{Welfare and FEPSDE}---the 144-component utility structure
that context modulates. Without knowing \emph{what} people value (FEPSDE), we cannot
assess how context shifts those values.

%------------------------------------------------------------------------------
\begin{cross-reference}
\textbf{Reading Path:}
\begin{itemize}
    \item For complete Ψ-dimension specification: Appendix WEN (CORE-WHEN)
    \item For temporal context: Appendix CTM (CONTEXT-TEMPORAL)
    \item For spatial context: Appendix CTS (CONTEXT-SPATIAL)
    \item For complexity economics: Appendix SHA (DOMAIN-COMPLEXITY)
    \item For welfare (next): Chapter 10 + Appendix WAT (CORE-WHAT)
    \item \textbf{For intervention optimization: Chapter 11 (EMERGE Algorithm) + Appendix IE (EIT-EMP)}
\end{itemize}
\end{cross-reference}

%------------------------------------------------------------------------------
\paragraph{Connection to Intervention Design (EIT-EMP).}
Context $\Psi$ is the universal modulator in the EMERGE algorithm (Chapter 11). Every
factor in the $K^*$ formula---$f_\gamma$, $f_S$, $f_C$, $f_L$---depends on context.
The endogenous nature of context ($\Psi_{t+1} = f(\Psi_t, a_t)$) means that interventions
change the context in which future interventions operate. This creates path dependence:
early intervention choices shape the feasibility and effectiveness of later ones. Chapter 21
(Dynamic Portfolio Evolution) addresses how to re-optimize $K^*$ as context evolves.

%%%%%%%%%%%%%%%%%%%%%%%%%%%%%%%%%%%%%%%%%%%%%%%%%%%%%%%%%%%%%%%%%%%%%%%%%%%%%%%
% END OF CHAPTER 9
%%%%%%%%%%%%%%%%%%%%%%%%%%%%%%%%%%%%%%%%%%%%%%%%%%%%%%%%%%%%%%%%%%%%%%%%%%%%%%%
